% tr-05 片段 b：原文 pp.38–40（PDF pp.39–41）
% 内容：5.1 续 —— (5.5)–(5.15)、(5.16)–(5.20)
% [原文 p.38]

其中 $y=\rho^2+\ii\tphi$。若比较 $J$ 的两个表达式，我们便得到原变量与对偶变量之间的关系。用 $y$ 与 $\phi$ 写出，它是 $(\partial_0+\partial_1)y=2\overline{\phi}(D_0+D_1)\phi$ 与 $(\partial_0-\partial_1)y=2\phi(D_0-D_1)\overline{\phi}$。又注意 $\Re(y)=|\phi|^2$。

对理论的 $(2,2)$ 超对称版本可以做同样的事，即对 $G=U(1)$、$V=\C(1)\ni\phi$、$W=0$、$\tW=-t\sigma$ 的 GLSM（也就是“$\CP^0$ 模型”）。把 $\phi$ 的相位对偶化之后，\footnote{存在一种明显保持 $(2,2)$ 超对称的做法 [8, 10]。}物质多重态 $(\phi,\psi,F)$ 变成一个超多重态 $(y,\eta,G)$，后者在超对称下的变换为

\begin{align}
    \delta y &= -\overline{\epsilon}_+\eta_- - \epsilon_-\eta_+, \notag\\
    \delta\eta_+ &= \ii\overline{\epsilon}_-(\partial_0+\partial_1)y-\overline{\epsilon}_+G,
        \qquad \delta\eta_- = \ii\epsilon_+(\partial_0-\partial_1)y+\epsilon_-G, \tag{5.5}\\
    \delta G &= \ii\epsilon_+(\partial_0-\partial_1)\eta_+ - \ii\overline{\epsilon}_-(\partial_0+\partial_1)\eta_-. \notag
\end{align}

与玻色情形一样，$\Im(y)$ 就是 $\phi$ 的相位 $\varphi$ 的 T-对偶像 $\tphi$。特别地，$y$ 在虚方向上周期化，

\begin{equation}
    y\equiv y+2\pi\ii.
    \tag{5.6}
\end{equation}

注意 $y$ 是 A-手征的：

\begin{equation}
    \overline{Q}_+y=Q_-y=0.
    \tag{5.7}
\end{equation}

正如我们在 3.2 节所见，对偶化把一个 B-手征多重态 $(\phi,\psi,F)$ 变成 A-手征多重态 $(y,\eta,G)$。因此它将是一个镜像对称。对偶化后理论的拉氏量由

\begin{equation}
    \widetilde{\cL}=Q_+Q_-\overline{Q}_+\overline{Q}_-\Bigl(-\frac{1}{2e^2}|\sigma|^2-\Re(y)\log\Re(y)\Bigr)
        +\Re\,Q_+\overline{Q}_-\bigl((y-t)\sigma\bigr).
    \tag{5.8}
\end{equation}
给出。注意对偶化后的理论具有扭超势

\begin{equation}
    \tW(\sigma,y)=(y-t)\sigma,
    \tag{5.9}
\end{equation}
这对应着如下事实：对偶的圆周变量 $\tphi=\Im(y)$ 扮演了动力学 theta 角的角色（见 (5.3)）。又注意 $K=\Re(y)\log\Re(y)$ 确实就是度规 (5.4) 的 Kähler 势。原变量与对偶变量之间的关系是

\begin{equation}
    \Re(y)=|\phi|^2,
    \tag{5.10}
\end{equation}
以及由它在超对称下推出的其他关系。例如

\begin{align}
    &\eta_+=2\overline{\phi}\psi_+,\qquad \eta_-=2\overline{\psi}_-\phi, \notag\\
    &(\partial_0+\partial_1)y=2\overline{\phi}(D_0+D_1)\phi-2\ii\overline{\psi}_+\psi_+, \tag{5.11}\\
    &(\partial_0-\partial_1)y=2\phi(D_0-D_1)\overline{\phi}+2\ii\overline{\psi}_-\psi_-. \notag
\end{align}
% [Tutor Note: p.38 (5.7) 是“$y$ 为 A-手征”的定义式：$\overline{Q}_+y=Q_-y=0$（注意第一个 $Q$ 带横杠、第二个不带；文字层把横杠整块丢掉）。可用 (5.5) 反查：$\delta y=-\overline{\epsilon}_+\eta_--\epsilon_-\eta_+$，故 $y$ 被 $\overline{Q}_+$ 与 $Q_-$ 湮灭正是“只含 $\eta_-$ 的 $\overline{Q}_+$ 分量、$\eta_+$ 的 $Q_-$ 分量”这一自由度计数的自然结果；与 3.2 节中“B-手征 $\to$ A-手征”的结论一致。]

% [Tutor Note: p.38 (5.8) 与 (5.9) 的自检：把 (5.8) 圆括号里的式子与 (5.18) 对照（取 $k=1$、$N=1$、$y_i^0\to y$）可见它就是 $-K$，只是 (5.8) 的规范动能项为 $\frac{1}{2e^2}$、而 (5.18) 印作 $\frac{1}{e_a^2}$（差一因子 2）；同页行内亦印作 $K=|\sigma|^2/(2e^2)+\cdots$、(5.23) 亦为 $\frac{1}{2e_a^2}$，故 (5.18) 疑漏印 $1/2$，详见 (5.18) 后的批注；而 $\tW$ 出现在 $\Re\,Q_+\overline{Q}_-\bigl((y-t)\sigma\bigr)$ 里，正是 (5.17) 的 $\Re\,Q_+\overline{Q}_-\tW$ 在 $\tW=(y-t)\sigma$ 时的样子。原文 p.38 的 “Written interms of” 缺一个空格，属排印笔误。]
% [原文 p.39]
注意当我们把费米子场置零时，这些关系就归结为上面得到的玻色情形。

关系 (5.10) 意味着 $\Re(y)$ 只能取非负值，即 $y$ 张成一个以 $\Re(y)=0$ 为端点的半圆柱，而度规 (5.4) 在该端点处发散。这个理论看起来也许很奇怪。然而，按照 FI 参数的跑动 $\zeta'=\zeta+\log(\mu'/\mu)$，我们对对偶变量作重正化 $y'=y+\log(\mu'/\mu)$，使得扭超势 (5.9) 在每个能标下都取相同的形式。特别地，截断标度 $\Lambda_0$ 处的变量 $y_0$ 与有限标度 $\mu$ 处的变量 $y$ 通过

\begin{equation}
    y_0=\log(\Lambda_0/\mu)+y.
    \tag{5.12}
\end{equation}
相联系。对偶化是在截断标度 $\Lambda_0$ 处进行的，因此界与度规分别是 $\Re(y_0)\ge0$ 与 $\dd\tilde{s}^2=|\dd y_0|^2/\bigl(2(y_0+\overline{y}_0)\bigr)$，也就是说

\begin{equation}
    \Re(y)\ge-\log(\Lambda_0/\mu),
    \tag{5.13}
\end{equation}
以及
\begin{equation}
    \dd\tilde{s}^2=\frac{|\dd y|^2}{2\bigl(2\log(\Lambda_0/\mu)+y+\overline{y}\bigr)}.
    \tag{5.14}
\end{equation}

当 $\Lambda_0/\mu\gg1$ 时，端点离得很远，度规近似为平坦圆柱的度规。

与重正化类似，$y$ 在（有异常的）$U(1)$ 轴转动下也作非平凡变换：

\begin{equation}
    \ee^{\ii\beta F_A}:\ y\longrightarrow y-2\ii\beta.
    \tag{5.15}
\end{equation}

这可以通过计算轴流与 (5.11) 中给出的 $(\partial_0\pm\partial_1)y$ 的算子乘积展开读出。又注意，theta 参数的轴平移（见 (4.22)；在目前情形下即 $\theta\to\theta+2\beta$）通过扭超势 (5.9) 被 (5.15) 有效地实现：$\tW=(y-\zeta+\ii\theta)\sigma$。还可以证明 $y$ 在矢量 $U(1)$ R-对称性下不变。

% [Tutor Note: p.39 (5.12)–(5.14) 的推导不必额外假设：由于 (5.9) 的 $\tW=(y-\zeta)\sigma$ 必须在每个能标下取相同形式，令 $\zeta'=\zeta+\log(\mu'/\mu)$ 就必须同时令 $y'=y+\log(\mu'/\mu)$，使 $y'-\zeta'=y-\zeta$；于是把截断标度 $\Lambda_0$ 处的变量写到标度 $\mu$ 处即得 (5.12)。(5.13)(5.14) 只是把 (5.12) 代入 $\Re(y_0)\ge0$ 与 (5.4)：$y_0+\overline{y}_0=2\log(\Lambda_0/\mu)+y+\overline{y}$。原文 p.39 “non-genative” 应作 “non-negative”（排印笔误，正文照译）。]

% [Tutor Note: p.39 式 (5.15) 的指数上是 $\beta F_A$（$F_A$ 为轴荷），文字层抽成空；核对方法：它必须与 (5.9) 中 $\theta\to\theta+2\beta$ 相容 —— $\ee^{\ii\beta F_A}$ 使 $y\to y-2\ii\beta$，而 $(y-\zeta+\ii\theta)\sigma$ 中的组合 $y+\ii\theta$ 在该变换下不变（因 $\theta\to\theta+2\beta$）。]

对具有多个 $U(1)$ 因子规范群和多个物质多重态的理论，对偶化的推广是直截了当的。让我们考虑一般模型：连通阿贝尔规范群且无超势，

\[
\begin{aligned}
    G&=U(1)^k,\\
    V&=\bigoplus_{i=1}^{N}\C(Q_i)\ni\phi=(\phi_1,\dots,\phi_N),\\
    W&=0,\\
    \tW&=-\sum_{a=1}^{k}t^a\sigma_a,
\end{aligned}
\]
% [原文 p.40]
其中我们再次假定 $G$ 在 $V$ 上的作用为忠实作用。相位的对偶化把 B-手征多重态 $(\phi_i,\psi_i,F_i)$ 变成 A-手征多重态 $(y_i,\eta_i,G_i)$，其中 $y_i$ 在虚部周期化，

\begin{equation}
    y_i\equiv y_i+2\pi\ii,
    \tag{5.16}
\end{equation}
而对偶化后理论的拉氏量为

\begin{equation}
    \widetilde{\cL}=Q_+Q_-\overline{Q}_+\overline{Q}_-(-K)+\Re\,Q_+\overline{Q}_-\tW,
    \tag{5.17}
\end{equation}

其中

\begin{equation}
    K=\sum_{a=1}^{k}\frac{1}{e_a^2}|\sigma_a|^2+\sum_{i=1}^{N}\Re(y_i^0)\log\Re(y_i^0),
    \tag{5.18}
\end{equation}
% [Tutor Note: 印刷 p.40 的 (5.18) 中规范动能项印作 $\frac{1}{e_a^2}|\sigma_a|^2$（1100 dpi 裁图确认分母无因子 2），这与印刷 p.38 的 (5.8) 的 $\frac{1}{2e^2}|\sigma|^2$、印刷 p.42 的 (5.23) 的 $\frac{1}{2e_a^2}|\sigma_a|^2$，以及同页行内 CP$^1$ 例的 $K=|\sigma|^2/(2e^2)+\cdots$ 均不一致；按上下文 (5.18) 应为 $\sum_a\frac{1}{2e_a^2}|\sigma_a|^2+\sum_i\Re(y_i^0)\log\Re(y_i^0)$。译文照原样排印（不擅改原文），此处仅作提示。]

\begin{equation}
    \tW=\sum_{a=1}^{k}\Bigl(\sum_{i=1}^{N}Q_i^a y_i-t^a\Bigr)\sigma_a.
    \tag{5.19}
\end{equation}

这里 $y_i^0$ 是截断标度 $\Lambda_0$ 处的对偶变量，它与物理标度 $\mu$ 处的 $y_i$ 通过 $y_i^0=\log(\Lambda_0/\mu)+y_i$ 相联系。变量 $y_i$ 在（可能具有异常的）轴 $U(1)$ R-转动下按 $y_i\to y_i-2\ii\beta$ 变换，并在矢量 $U(1)$ R-对称性下不变。

让我们稍作停顿，看看在 $\CP^1$ 模型中会得到什么。此时 $K=|\sigma|^2/(2e^2)+\Re(y_1^0)\log\Re(y_1^0)+\Re(y_2^0)\log\Re(y_2^0)$，而 $\tW=(y_1+y_2-t)\sigma$。回忆一下，当 $e\gg|\Lambda|=\ee^{-\zeta/2}\mu$ 时（例如当 $e\to\infty$ 时），GLSM 可以看作 $\CP^1$ sigma 模型。我们可以先积掉规范多重态，得到一个约束 $y_1+y_2=t$，于是剩下的是由 $y=y_1$ 参数化的靶空间的 sigma 模型，其度规为

\begin{equation}
    \dd\tilde{s}^2=\frac{|\dd y|^2}{2\bigl(2\log(\Lambda_0/\mu)+y+\overline{y}\bigr)}
        +\frac{|\dd y|^2}{2\bigl(2\log(\Lambda_0/\mu)+2\zeta-y-\overline{y}\bigr)}.
    \tag{5.20}
\end{equation}

注意 $y^0=y+\log(\Lambda_0/\mu)$ 的实部被界定为 $0\le\Re(y^0)\le\zeta_0$，其中 $\zeta_0=\zeta+2\log(\Lambda_0/\mu)$ 是截断标度 $\Lambda_0$ 处的 FI 参数。这实质上与图 2 所描绘的是同一个几何，而我们知道它作为 $\CP^1$ 的对偶会有一些严重的问题。因此，上面的结果不可能是事情的全貌。

% [Tutor Note: p.40 (5.20) 的自检：由约束 $y_1+y_2=t$ 立刻得 $\Re(y_1^0)+\Re(y_2^0)=2\log(\Lambda_0/\mu)+\Re(t)=\zeta_0$，所以两个对数项的分母分别是 $2\Re(y^0)$ 与 $2(\zeta_0-\Re(y^0))$，这正是 (5.20) 的两项；由此立见 $\Re(y^0)$ 被夹在 $0$ 与 $\zeta_0$ 之间（正文那句话），也说明当 $\Re(y^0)\to0$ 或 $\to\zeta_0$ 时度规发散 —— 这就是“图 2 的几何有严重问题”的来源。原文 p.40 “with taget space” 应作 “with target space”（排印笔误，正文照译）。]

类似地，在一般模型中，积掉规范多重态之后，我们得到环面纤维化 $\mu_F:X_{\zeta_0}\to P_{\zeta_0}$ 的朴素对偶所对应的 sigma 模型，其中 $\zeta_0$ 是截断标度处的 FI 参数。满足界 $\Re(y_i^0)\ge0$ 与约束 $\sum_i Q_i^a\Re(y_i^0)=\zeta_0^a$ 的诸 $\Re(y_i^0)$ 通过 $(\Re(y_i^0))\mapsto\bigl(\sum_i f_i^\alpha\Re(y_i^0)\bigr)$ 同构地投影到矩多面体 $P_{\zeta_0}$ 上，而满足 $\sum_i Q_i^a\Im(y_i^0)\equiv-\theta^a\ (\mathrm{mod}\ 2\pi\Z)$ 的诸 $\Im(y_i^0)$ 则参数化对偶纤维。这与 $\CP^1$ 模型中的问题相同，因而同样不可能是事情的全貌。

% [Tutor Note: p.40 末尾的 $\zeta_0^a$ 是上标 $a$、下标 $0$（文字层抽成 “ζa 0”，极易抄成别的形状），与同段的 $X_{\zeta_0},P_{\zeta_0}$ 共用 $\zeta_0$；$f_i^\alpha$ 是投影矩阵（文字层抽成 “f α i”）。判据：投影后落在 $P_{\zeta_0}$ 的 $\alpha$ 坐标里，故 $\alpha$ 指标必须与多面体的维数配对，而 $i$ 是物质场的指标。]
